\documentclass[10pt]{mypackage}

\usepackage[uprightgreeks,light]{kpfonts}
\renewcommand{\coloneq}{\coloneqq}
%\renewcommand{\mathbb}{\mathds}
\usepackage{homework}
%\usepackage{notes}

%\usepackage[ backend=bibtex, style = alphabetic, sorting=ynt ]{biblatex}
%\addbibresource{  }

\usepackage{parskip}

\fancyhf{}
\fancyhead[R]{Avinash Iyer}
\fancyhead[L]{Algebra II: Homework 10}
\fancyfoot[C]{\thepage}

\setcounter{secnumdepth}{0}

\begin{document}
\RaggedRight
\begin{problem}[Problem 1]
  Prove the interesting part of Corollary 18.7 from the Spring 2010 notes: if $K/L/F$ is a tower of algebraic extensions, and $K/L$ and $L/F$ are separable, then $K/F$ is separable.
\end{problem}
\begin{solution}
  Suppose $K/L$ and $L/F$ are separable, with $K/F$ a finite extension. Then, since $\left[ K:F \right] = \left[ K:L \right]\left[ L:F \right]$, we must have that both $K/L$ and $L/F$ are finite extensions. It follows that there are towers of separable simple extensions $L_0 =L \subseteq L_1\subseteq\cdots\subseteq L_m = K$, and $F_0 = F \subseteq F_1\subseteq\cdots\subseteq F_n = L$. Combined, this yields a tower of separable simple extensions 
  \begin{align*}
    F_0 = F \subseteq F_1\subseteq \cdots \subseteq F_n\subseteq L_1\subseteq\cdots\subseteq L_m = K,
  \end{align*}
  so that $K/F$ is separable.

  If $K/F$ is not a finite extension, we let $\alpha\in K$ be any element, and choose some polynomial $p\in K[x]$ given by
  \begin{align*}
    p(x) &= a_0 + a_1 x + \cdots + a_{m-1}x^{m-1} + x^{m}
  \end{align*}
  that has $p(\alpha) = 0$. We then have that $M = F\left( a_0,\dots,a_{m-1} \right)/F$ is a finite extension. Since we have either $M\subseteq L$ or $L\subseteq M$, we have either $L/M/F$ is a tower of separable extensions or $K/M/L$, following from the fact that $L/F$ and $K/L$ respectively are separable. Let $M' = M\left( \alpha \right)$, so that $M'/F$ is a finite extension. If $M'\subseteq L$, then any element of $M'$ is an element of $L$, so the minimal polynomial of any element of $M'$ is separable over $F$, meaning $M'/F$ is separable. If $L\subseteq M'$, then the tower $M'/L/F$ has $\left[ M':F \right]$ finite, so that $\left[ M':L \right]$ and $\left[ L:F \right]$ are finite separable extensions, reducing to the previous case of $K/F$ as finite.

  Therefore, any element $\alpha\in K$ is contained in a separable extension over $F$, so $K$ is separable over $F$.
\end{solution}
\begin{problem}[Problem 2]
Let $p_1,\dots,p_n$ be distinct primes, and let $\alpha_i = \sqrt[p_i]{p_i}$, $\alpha = \sum_{i=1}^{n}\alpha_i$. Prove that $\Q\left( \alpha_1,\dots,\alpha_n \right) = \Q\left( \alpha \right)$.
\end{problem}
\begin{problem}[Problem 3]
  Let $p$ be a prime, $K = \F_p\left( s,t \right)$, the field of rational functions over $\F_p$ in two variables, and let $F = \F_p\left( s^p,t^p \right)$. Prove that the extension $K/F$ cannot be generated by a single element.
\end{problem}
\begin{problem}[Problem 4]
  Let $f(x)\in \Q[x]$ be an irreducible polynomial of degree $n$, and let $K\subseteq \C$ be the splitting field for $f(x)$ over $\Q$. Label the roots of $f(x)$ by $1,\dots,n$ in some order, and let $\iota\colon \operatorname{Gal}\left( K/\Q \right)\rightarrow S_n$ be the associated embedding.
  \begin{enumerate}[(a)]
    \item Prove that $\left\vert \operatorname{Gal}\left( K/\Q \right) \right\vert$ is a multiple of $n$ and divides $n!$.
    \item Assume $f(x)$ has at least one non-real root. Prove that complex conjugation is an element of $ \operatorname{Gal}\left( K/\Q \right) $ of order $2$.
    \item Assume that $f(x)$ has precisely two non-real roots. Prove that the image of the complex conjugation under the embedding $\iota$ is a transposition.
    \item Suppose that $n = \operatorname{deg}\left( f \right)$ is prime, and assume that $f(x)$ has exactly two non-real roots. Prove that $\operatorname{Gal}\left( K/\Q \right)$ is isomorphic to $S_n$.
  \end{enumerate}
\end{problem}
\begin{solution}\hfill
  \begin{enumerate}[(a)]
    \item We have that $\left\vert \operatorname{Gal}\left( K/\Q \right) \right\vert\subseteq S_n$ is a subgroup, so Lagrange's theorem immediately gives that $\left\vert \operatorname{Gal}\left( K/\Q \right) \right\vert$ divides $n!$.

      Label the roots of $f$ by $\set{\alpha_1,\dots,\alpha_n}$. Note that the Galois group acts transitively on this set, so by orbit--stabilizer, we have
      \begin{align*}
        n &= \left\vert O_{\alpha_1} \right\vert\\
          &= \left\vert G \right\vert/\left\vert G_{\alpha_1} \right\vert,
      \end{align*}
      where we have $G = \operatorname{Gal}\left( K/\Q \right)$, so that $n$ divides $\left\vert \operatorname{Gal}\left( K/\Q \right) \right\vert$.
    \item 
  \end{enumerate}
\end{solution}
\end{document}
